\scanpage{4}
% The typed line at the head of the sheet is keyed by a pen loop to the first case below.

et en vertu de $\overset{\handwritten{\sim}}{c}$), on a $u\handwritten{\in} \overset{\handwritten{\sim}}{M}$ sss pour tout $i$, $F\times_G X_i\to X_i$ est dans $\overset{\handwritten{\sim}}{M}$. Cela nous ramène au cas où $G$ est de la forme $X\handwritten{\in} S$. Je dis qu'alors $u$ est dans $\overset{\handwritten{\sim}}{M}$ sss il existe une famille couvrante $u_i:X_i\to F$ avec \deleted{$u_i\unsure{\in} \overset{\handwritten{\sim}}{M}$, $uu_i\unsure{\in} M$,} $X_i\handwritten{\in} S$, $uu_i\handwritten{\in} M$, \deleted{\unsure{\&test|et est}} et les $X_i\times_F X_j\xrightarrow{\mathrm{pr}_1}X_i$ dans $M$ (NB cela a un sens, car l'hypothèse faite sur $u$ implique que $X_i\times_F X_j$ est représentable, comme image inverse \inserted{dans $X_i\times_X X_j$} de la diagonale de $F$ sur $X$)\deleted{dans} La condition est nécessaire:\deleted{\unsure{=|:}}on trouve les $u_i$ grâce à $\overset{\handwritten{\sim}}{f}$), alors par $\overset{\handwritten{\sim}}{a}$) la projection $X_i\times_F X_j\to X_i$ est dans $\overset{\handwritten{\sim}}{M}$, donc dans $M$ par $\overset{\handwritten{\sim}}{e}$)\handwritten{\unsure{6|$\ell$}}\inserted{\handwritten{;}les $uu_i$ sont dans $M$ par $\overset{\handwritten{\sim}}{b}$) et $\overset{\handwritten{\sim}}{e}$).} La condition est suffisant\unsure{e|é} grâce à $\overset{\handwritten{\sim}}{e}$) et $\overset{\handwritten{\sim}}{d}$).

2º) Cas général. Procédant comme plus haut, on estbramené au cas où $G$ est un $X\handwritten{\in} S$, et dans ce cas, on voit de même que $u\handwritten{\in} \overset{\handwritten{\sim}}{M}$ sss il existe une famille couvrante de flèches $u_i:X_i\handwritten{\to} F$, avec $uu_i\handwritten{\in} M$, et les $Z_{ij}=X_i\times_F X_j\xrightarrow{\mathrm{pr}_1}X_i$ dans $\overset{\handwritten{\sim}}{M}$. Or ici $Z_{ij}$ est un sous-\deleted{\unsure{foncteur}}faisceau\unsure{.} du faisceau représentable $Z'_{ij}=X_i\times_X X_j$, donc \deleted{sa}son \deleted{diagonale}\inserted{morphisme diagonal} sur $X_i$ est image inverse du morphisme diagonal pour $Z'_{ij}$ sur $X_i$, qui est représentable car le carré fibré de $Z'_{ij}$ sur $X_i$ l'est, donc delui de $Z_{ij}$ sur $X_i$ l'est aussi. On est donc ramené au cas 1º).

Pour achever de prouver la proposition, il reste à montrer que l'ensemble de flèches $\overset{\handwritten{\sim}}{M}$ construit par les critères précédents satisfait bien aux conditions $\overset{\handwritten{\sim}}{a}$) à $\overset{\handwritten{\sim}}{f}$), ce qu'on laisse au lecteur le soin de faire.

\uline{Définition} On dit qu'un objet $F$ de $\overset{\handwritten{\sim}}{S}$ peut \deleted{\unsure{a}}s\uline{'obtenir par recollement d'objets de} $S$ \uline{à l'aide de morphismes dans} $M$, si on peut trouver une famille couvrante $u_i:X_i\to F$, avec $X_i\handwritten{\in} S$ et $u_i\handwritten{\in} \overset{\handwritten{\sim}}{M}$. On désignera par $ {\overset{\handwritten{\sim}}{S}}_M$ la sous-catégorie pleine de $\overset{\handwritten{\sim}}{S}$ ayant comme objets ces $F$.

\deleted{\unsure{n}}On a alors \deleted{\unsure{le}}la propriété de stabilité suivante:

\deleted{\uline{Proposition} \quad 1) Soit $F\quad\mathrm{Ob}\,\overset{\handwritten{\sim}}{S}$, et $F_{i}\quad F$ une famille couvrante}

% archive sheet number: 3
